\documentclass[10pt,a4paper]{article}

\usepackage[margin=1.55cm]{geometry}
\usepackage{graphicx}
\usepackage{float}
\usepackage{booktabs}
\usepackage{amsmath}
\usepackage{xcolor}
\usepackage{hyperref}
\usepackage{caption}
\usepackage{subcaption}
\usepackage{enumitem}
\usepackage{listings}
\usepackage{setspace}
\usepackage{titlesec}

\setstretch{0.96}
\setlength{\parskip}{2pt}
\setlength{\parindent}{0pt}
\setlist[itemize]{noitemsep, topsep=2pt, leftmargin=14pt}
\setlist[enumerate]{noitemsep, topsep=2pt, leftmargin=16pt}

\titlespacing*{\section}{0pt}{7pt}{4pt}
\titlespacing*{\subsection}{0pt}{5pt}{2pt}

\captionsetup{font=footnotesize, labelfont=bf}

\hypersetup{
	colorlinks=true,
	linkcolor=blue,
	urlcolor=blue,
	citecolor=blue
}

\newcommand{\figdir}{results/figures/report/}

\newcommand{\safeincludegraphics}[2][0.9\linewidth]{%
	\IfFileExists{\figdir#2}{%
		\includegraphics[width=#1]{\figdir#2}%
	}{%
		\fbox{%
			\parbox[c][3cm][c]{0.85\linewidth}{%
				\centering
				Missing figure:\\
				\texttt{\detokenize{#2}}%
			}%
		}%
	}%
}

\lstset{
	basicstyle=\ttfamily\scriptsize,
	breaklines=true,
	frame=single,
	columns=fullflexible,
	backgroundcolor=\color{gray!5}
}

\begin{document}
	
	\begin{titlepage}
		\centering
		\vspace*{1.6cm}
		
		{\Large Budapest University of Technology and Economics\\}
		{\large Faculty of Electrical Engineering and Informatics\\[1.2cm]}
		
		{\Huge \textbf{xAI-Based Alarm Classification for High-Recall Traffic Incident Detection
			}\\[0.4cm]}
		{\Large Project Laboratory Report\\[1.1cm]}
		
		\begin{tabular}{rl}
			\textbf{Student:}        & Arbenite Canaj \\
			\textbf{Program:}        & BSc Computer Engineering \\
			\textbf{Course:}         & Project Laboratory \\
			\textbf{Course code:}    & BMEVIHIAL04 \\
			\textbf{Supervisor:}     & Prof. Mohammad Bawaneh \\
		\end{tabular}
		
		\vfill
		{\large \today}
	\end{titlepage}
	
	\section*{Abstract}
	\addcontentsline{toc}{section}{Abstract}
	
	\begin{center}
		\begin{minipage}{0.78\textwidth}
			
		Many AI models behave like black boxes, meaning that while they are capable of making accurate predictions, it is sometimes unknown why they came to a specific conclusion. The lack of transparency becomes a major issue as more and more non-experts rely on these systems, particularly in high-stakes situations where ambiguous or unreliable results can cause individuals to trust the system for the wrong reasons, with serious consequences. In response, this research uses Explainable Artificial Intelligence (XAI) to detect traffic incidents, where operators have to make decisions in real time under uncertainty.	
			
		\vspace{6pt}
		
		Automatic Incident Detection (AID) systems output a binary alarm indicating whether a traffic anomaly has been detected. Without further context, operators have limited support for judging whether an alarm deserves immediate attention. By revealing the temporal traffic evidence underlying model conclusions and providing operators with supporting data for alarm classification and evaluation, this study explores whether XAI techniques may close this gap.
		
		\vspace{6pt}
		This research uses prepared and smoothed traffic incident-window data from 175 sensor streams, building on the referenced work\cite{abujudom2025}. The earlier work provides the data and protocol framework, but this project focuses on a distinctive contribution, an XAI-based alarm support layer. A comparison is made between four families of anomaly-detection models: Random Forest, XGBoost, Multi-Layer Perceptron (MLP), and Long Short-Term Memory network (LSTM). The primary input features used by the models are occupancy and speed.
		
		\vspace{6pt}
		Since the input space is limited to speed and occupancy, the main contribution is not only about identifying which input feature matters most.  XAI is used to describe the quality and structure of alarm evidence. Alarm-positive timestamps are categorized into alarm episodes, and score margin, model agreement, recent attribution share, explanation concentration, SHAP/LIME agreement, and cross-model explanation agreement can all be used to assess each episode.
		
		\vspace{6pt}
		XAI is not used to automatically detect false positives in real time or to demonstrate the actual physical reason of an alarm. Rather, it illustrates the recent speed and occupancy context that caused an alarm in order to understand model behavior. From an operator-support standpoint, this enables the alarm support layer to distinguish between stronger, borderline, and less stable alarms; nonetheless, the human operator retains the final say.
		
			
			
		\end{minipage}
	\end{center}
	
	\newpage
	
	\section{Introduction}
	
	Traffic incidents can quickly affect the performance of a road network. Even a relatively short disruption may lead to congestion, longer travel times, and less reliable traffic flow. For this reason, traffic management centers depend on monitoring systems that can detect unusual traffic conditions early enough to support a timely response.
	
	Automatic Incident Detection (AID) has been studied for many years, and machine learning has made it possible to recognize abnormal traffic patterns more effectively. Still, good detection performance is not enough on its own. In practice, an operator also needs to understand how serious an alarm appears to be and whether it is supported by clear traffic evidence.
	
	This project investigates whether additional information derived from Explainable Artificial Intelligence (XAI) can improve the interpretation of incident alarms. Rather than focusing solely on whether an anomaly is detected, the work examines how model explanations and temporal traffic context can be used to characterize the supporting evidence behind an alarm. The goal is to provide a richer basis for alarm assessment and prioritization in operational traffic management environments.
	
	The main research question is:
	
	\begin{quote}
		How can XAI methods and temporal traffic context be combined to provide traffic operators with useful supporting information for judging incident alarms?
	\end{quote}
	
	The project investigates the following sub-questions:
	
	\begin{itemize}
		\item How do different model families explain traffic anomaly alarms?
		\item Do SHAP and LIME identify similar temporal evidence for the same alarm?
		\item Can XAI-derived evidence help distinguish high-priority alarms from borderline alarms?
		\item How can model explanations be translated into operator-facing alarm cards?
	\end{itemize}
	
	\section{Dataset and Experimental Setup}
	
	\subsection{Data Source}
	
	This project uses prepared traffic incident data derived from the framework of the previous thesis \cite{abujudom2025}. The source files are individual incident CSV streams, not one unified raw traffic dataset. Each stream corresponds approximately to one incident timeline.
	
	The input files are stored as raw and smoothed incident files. The consolidated processed dataset contains 175 selected incident streams. Each stream contains:
	
	\begin{itemize}
		\item 240 hours of baseline training data,
		\item 12 hours of test data,
		\item 5-minute sampling intervals,
		\item smoothed speed and smoothed occupancy features,
		\item incident labels for the target incident period inside the test window.
	\end{itemize}
	
	Since the data is sampled every 5 minutes, each stream contains 2880 training rows and 144 test rows, for a total of 3024 rows per stream. Across 175 streams, the consolidated dataset contains 529,200 rows.
	
	\begin{table}[H]
		\centering
		\caption{Dataset overview.}
		\begin{tabular}{lr}
			\toprule
			Property & Value \\
			\midrule
			Incident streams & 175 \\
			Training window per stream & 240 hours \\
			Test window per stream & 12 hours \\
			Sampling interval & 5 minutes \\
			Train rows per stream & 2880 \\
			Test rows per stream & 144 \\
			Total rows & 529,200 \\
			Input features & speed\_smoothed, occ\_smoothed \\
			EWMA smoothing alpha & 0.15 \\
			\bottomrule
		\end{tabular}
	\end{table}
	
	Only smoothed speed and smoothed occupancy are used as model inputs. Flow is not used in this project. Therefore, all XAI explanations are necessarily limited to speed and occupancy behavior over time.
	
	\begin{figure}[H]
		\centering
		\safeincludegraphics[0.95\linewidth]{project_pipeline_overview.png}
		\caption{End-to-end project workflow from prepared incident files to operator-facing alarm classification.}
		\label{fig:project_pipeline}
	\end{figure}
	
	
	\begin{figure}[H]
		\centering
		\safeincludegraphics[0.9\linewidth]{incident_window_protocol.png}
		\caption{Incident-window protocol. Each stream contains a 240-hour incident-free baseline window followed by a 12-hour test window containing the labeled target incident period.}
		\label{fig:incident_window_protocol}
	\end{figure}
	

	
	\subsection{Splitting Protocol}
	
	Two stages of splitting are used in the project.
	
	Initially, each stream is internally divided into a 12-hour test window and a 240-hour baseline training window. The baseline window is used to learn typical traffic behavior and is incident-free. 
	Second, there are five folds of 35 streams each, totaling 175 streams. To prevent leaks, entire streams are kept together.
	Four folds are used for training and calibration during each cross-validation cycle, with one fold being kept out for evaluation.
	
	\begin{table}[H]
		\centering
		\caption{Cross-validation design.}
		\begin{tabular}{lr}
			\toprule
			Property & Value \\
			\midrule
			Number of folds & 5 \\
			Streams per fold & 35 \\
			Evaluation unit & Whole streams \\
			Train data & Incident-free baseline windows \\
			Test data & 12-hour incident windows \\
			Leakage prevention & Streams kept intact \\
			\bottomrule
		\end{tabular}
	\end{table}
	
	\begin{figure}[H]
		\centering
		\safeincludegraphics[0.8\linewidth]{stream_level_cv_diagram.png}
		\caption{Five-fold stream-level cross-validation. The 175 streams are split into five groups of 35 streams, and each group is held out once. }
		\label{fig:stream_cv}
	\end{figure}
	
	
	\section{Model Setup}
	
	Four model families are compared: Random Forest, XGBoost, MLP, and LSTM. The models are not trained as standard supervised incident classifiers. Instead, they are used as anomaly detectors based on baseline prediction or residual scoring. They learn normal traffic behavior from incident-free baseline data. At test time, a large difference between predicted and observed traffic behavior becomes an anomaly score.
	
	\begin{table}[H]
		\centering
		\caption{Models and explanation methods.}
		\begin{tabular}{lll}
			\toprule
			Model & Model family & XAI method \\
			\midrule
			Random Forest & Tree ensemble & TreeSHAP, LIME \\
			XGBoost & Boosted tree ensemble & TreeSHAP, LIME \\
			MLP & Neural network & DeepSHAP-style attribution, LIME \\
			LSTM & Sequence neural network & Sequence DeepSHAP-style attribution \\
			\bottomrule
		\end{tabular}
	\end{table}
	
	The Random Forest, XGBoost, and MLP models use flattened temporal lookback windows. With a lookback of 12 timestamps and two input variables, the explainable input contains lagged features such as:
	
	\begin{itemize}
		\item \texttt{speed\_smoothed\_t\_minus\_1}: speed 5 minutes before the alarm,
		\item \texttt{speed\_smoothed\_t\_minus\_2}: speed 10 minutes before the alarm,
		\item \texttt{occ\_smoothed\_t\_minus\_1}: occupancy 5 minutes before the alarm,
		\item \texttt{occ\_smoothed\_t\_minus\_2}: occupancy 10 minutes before the alarm.
	\end{itemize}
	
	LIME was not applied to the LSTM in this project because the standard LIME setup is more suitable for flattened tabular inputs. LIME explains a prediction by creating perturbed versions of the input and observing how the model output changes. For the Random Forest, XGBoost, and MLP models, this is straightforward because the 60-minute lookback window is flattened into input columns. In contrast, the LSTM receives the input as an ordered sequence of time steps. Randomly perturbing individual values in this sequence can break the temporal structure of the traffic pattern and may create unrealistic sequence inputs. Therefore, sequence DeepSHAP-style attribution was used for the LSTM, since it preserves the time-step structure and provides attribution across both time and feature dimensions.
	
	\section{Evaluation Metrics}
	
	Two evaluation views are used.
	
	The first view is timestamp-level evaluation. This evaluates individual 5-minute rows. Precision measures how many raised alarms were correct incident timestamps. Recall measures how many incident timestamps were detected. F1-score summarizes the balance between precision and recall.
	
	The second view is incident-level evaluation, following the style of the referenced thesis. Detection Rate measures whether each incident was detected at least once. 
	
	False Alarm Rate is calculated as:
	
	\[
	FAR = \frac{\text{false alarm timestamps}}{\text{normal test timestamps}}
	\]
	
	A false alarm timestamp is a 5-minute test row outside the labeled incident period where the model still raises an alarm. 
	
	Mean Time to Detection measures the average delay between the start of an incident and the first alarm.
	
	
	\begin{figure}[H]
		\centering
		\safeincludegraphics[0.78\linewidth]{timestamp_level_model_comparison.png}
		\caption{Timestamp-level precision, recall, and F1-score for the final model outputs.}
		\label{fig:timestamp_level_comparison}
	\end{figure}
	
	
	\begin{figure}[H]
		\centering
		\safeincludegraphics[0.92\linewidth]{incident_level_model_comparison.png}
		\caption{Incident-level model comparison using Detection Rate, False Alarm Rate, and Mean Time to Detection. }
		\label{fig:incident_level_comparison}
	\end{figure}
	
	Figure~\ref{fig:incident_level_comparison} is closer to the operator's perspective because an operator cares whether the incident was detected and how quickly the first alarm appeared.
	
	The final incident-level values show that all models detect most incidents, while still producing false alarm timestamps due to the high-recall operating point. For example, Random Forest has \(1240\) false alarm timestamps over \(21276\) normal test timestamps, giving \(1240/21276 = 0.0583\). The final values are summarized in Table~\ref{tab:incident_values}.
	
	\begin{table}[H]
		\centering
		\caption{Incident-level evaluation values.}
		\label{tab:incident_values}
		\begin{tabular}{lrrr}
			\toprule
			Model & Detection Rate & False Alarm Rate & MTTD (min) \\
			\midrule
			Random Forest & 0.994 & 0.0583 & 0.686 \\
			XGBoost & 1.000 & 0.0579 & 1.086 \\
			MLP & 0.983 & 0.0580 & 0.724 \\
			LSTM & 0.994 & 0.0499 & 0.544 \\
			\bottomrule
		\end{tabular}
	\end{table}
	
	\begin{figure}[H]
		\centering
		\safeincludegraphics[0.75\linewidth]{fold_performance_stability.png}
		\caption{Per-fold timestamp recall stability. }
		\label{fig:fold_stability}
	\end{figure}
	
	Figure~\ref{fig:fold_stability} shows that recall varies across folds because each held-out group contains different incident streams. Random Forest and XGBoost are relatively stable, while MLP and LSTM show more fold-to-fold variation. This supports using stream-level cross-validation instead of relying on one split, because incident detection difficulty depends on the specific traffic cases being evaluated.
	
	\section{Alarm Episode Classification Layer}
	
	\subsection{Timestamp Predictions to Alarm Episodes}
	
	The models produce predictions every five minutes. However, a sequence of consecutive alarm-positive timestamps should not be treated as separate operational alarms. Therefore, consecutive alarm-positive timestamps are grouped into alarm episodes.
	
	For each alarm episode, one representative row is selected. This row is an actual timestamp row from the dataset, usually the strongest or most informative point inside the episode, such as the row with the highest anomaly score. The representative row is used to summarize and explain the alarm episode.
	
	\begin{figure}[H]
		\centering
		\safeincludegraphics[0.95\linewidth]{example_alarm_timeline.png}
		\caption{Example alarm timeline showing smoothed speed, smoothed occupancy, labeled incident period, and alarm-positive timestamps from the models.}
		\label{fig:example_alarm_timeline}
	\end{figure}

	
	\begin{figure}[H]
		\centering
		\safeincludegraphics[0.9\linewidth]{timestamp_to_alarm_episode_grouping.png}
		\caption{Grouping timestamp-level alarm predictions into one alarm episode. The representative row summarizes the episode for XAI and operator review.}
		\label{fig:timestamp_to_episode}
	\end{figure}
	
	
	\section{Explainable AI Methodology}
	
	\subsection{Role of XAI in This Project}
	
	XAI is used as an alarm evidence-quality layer. It does not determine the actual physical cause of an alarm or demonstrate that an incident occurred. Rather, it explains the behavior of the model by highlighting the aspects of the recent speed and occupancy history that were most important in determining the anomaly score.
	
	
	Two different kinds of evidence are combined in the alarm classification layer:
	
	\begin{table}[H]
		\centering
		\caption{Non-XAI and XAI evidence used for alarm classification.}
		\begin{tabular}{lp{9cm}}
			\toprule
			Evidence type & Examples \\
			\midrule
			Non-XAI alarm strength & anomaly score, threshold, score margin, alarm duration, number of models alarming \\
			XAI evidence structure & top features, recent attribution share, top-3 concentration, SHAP/LIME overlap, cross-model explanation agreement \\
			\bottomrule
		\end{tabular}
	\end{table}
	
	This separation is important. Score margin and model agreement describe how strongly the system alarmed. XAI describes what temporal traffic evidence supported that alarm.
	
	\subsection{SHAP}
	
	SHAP \cite{shap} explains a model output by assigning contribution values to input features. A simplified local explanation can be written as:
	
	\[
	f(x) = \phi_0 + \sum_{j=1}^{M} \phi_j
	\]
	
	Here, \(f(x)\) is the model output for one alarm row, \(\phi_0\) is the reference output, and \(\phi_j\) is the contribution of input feature \(j\). In this project, an input feature is a lagged speed or occupancy value, such as speed 5 minutes before the alarm.
	
	For tree-based models, TreeSHAP is used. For the MLP, a DeepSHAP-style attribution method is used. For the LSTM, sequence attribution is used over the time-by-feature input.
	
	\subsection{LIME}
	
	LIME \cite{lime} explains a prediction by generating perturbed versions of an input and fitting a simple local surrogate model around that prediction. A simplified LIME surrogate can be written as:
	
	\[
	g(z) = \beta_0 + \sum_{j=1}^{M} \beta_j z_j
	\]
	
	Here, \(g(z)\) is the local interpretable model fitted near one prediction, and \(\beta_j\) describes the local influence of feature \(j\). In this project, LIME is applied to Random Forest, XGBoost, and MLP. LIME is not currently used for the LSTM because the LSTM input is sequential and the current implementation does not support LIME for that architecture.
	
	SHAP and LIME values are not compared directly because they use different scales and assumptions. Instead, agreement is measured using top-feature overlap, and temporal emphasis.
	
	
	\subsection{Temporal Evidence Metrics}
	
	Because the project uses only speed and occupancy, the explanation analysis should not stop at the level of speed versus occupancy. The more meaningful analysis is temporal. Each explanation is summarized using:
	
	\begin{itemize}
		\item exact top features,
		\item speed versus occupancy attribution share,
		\item recent 0--15 minute attribution share,
		\item mid 15--30 minute attribution share,
		\item older 30--60 minute attribution share,
		\item top-3 attribution concentration,
		\item SHAP/LIME agreement where available,
		
	\end{itemize}
	
	
	\section{Cross-Fold XAI Evidence Analysis}
	
	The final XAI analysis goes beyond isolated examples. For each fold, alarm episodes were first identified, and one representative alarm row was selected from each episode. XAI evidence metrics were then calculated for these representative rows across all five folds. Therefore, the following figure is not a single alarm explanation; it summarizes repeated explanation patterns across the full alarm set.
	
	\begin{figure}[H]
		\centering
		\safeincludegraphics[0.92\linewidth]{shap_style_summary_beeswarm_all_episodes.png}
		\caption{SHAP-style feature influence summary across all alarm episode representative rows.}
		\label{fig:shap_beeswarm}
	\end{figure}

Figure~\ref{fig:shap_beeswarm} shows that the strongest repeated explanation pattern is concentrated in the most recent traffic measurements. The largest horizontal spread appears for speed and occupancy at 5 minutes before the representative alarm row, meaning these features most often carried the largest explanation weight. This supports the conclusion that many alarm explanations are driven by very recent traffic changes rather than by older traffic history.

For speed, low recent values are the important congestion-like signal, because a sudden speed drop is typical during abnormal traffic conditions. For occupancy, high recent values are the important congestion-like signal, because high occupancy means the road segment is more crowded. Therefore, the plot shows that recent low speed and recent high occupancy are the clearest recurring evidence patterns behind the alarms. Features from 20--60 minutes before the alarm are clustered much closer to zero, so they generally contribute less to the explanation.
	
	\begin{figure}[H]
		\centering
		\safeincludegraphics[0.82\linewidth]{temporal_attribution_heatmap_all_episodes.png}
		\caption{Temporal attribution heatmap across all alarm episodes. The figure shows whether model explanations are mainly based on recent traffic evidence, mid-range evidence, or older traffic history.}
		\label{fig:temporal_xai_heatmap}
	\end{figure}
	
	Figure~\ref{fig:temporal_xai_heatmap} is central to the project because the operational question is temporal: whether alarms are supported by recent local traffic context or by older history.
	
	The temporal evidence results show clear differences between model-method combinations. LSTM sequence explanations are highly recent-focused, with about \(97.7\%\) of attribution in the 0--15 minute window. TreeSHAP for Random Forest and XGBoost is also strongly recent-focused, while MLP explanations are more spread across the 60-minute lookback. XGBoost with LIME-style explanations shows more older-context evidence than TreeSHAP. 
	
	\begin{table}[H]
		\centering
		\caption{Mean XAI evidence metrics by model and explanation method.}
		\label{tab:xai_evidence_summary}
		\resizebox{\linewidth}{!}{%
			\begin{tabular}{llrrrrr}
				\toprule
				Model & Method & Speed share & Occupancy share & Recent 0--15 min & Mid 15--30 min & Older 30--60 min \\
				\midrule
				Random Forest & TreeSHAP & 0.502 & 0.498 & 0.987 & 0.009 & 0.004 \\
				Random Forest & LIME & 0.516 & 0.484 & 0.712 & 0.093 & 0.195 \\
				XGBoost & TreeSHAP & 0.509 & 0.491 & 0.967 & 0.021 & 0.012 \\
				XGBoost & LIME & 0.514 & 0.486 & 0.539 & 0.155 & 0.305 \\
				MLP & DeepSHAP-style & 0.504 & 0.496 & 0.771 & 0.106 & 0.123 \\
				MLP & LIME & 0.524 & 0.476 & 0.734 & 0.108 & 0.158 \\
				LSTM & Sequence DeepSHAP & 0.548 & 0.452 & 0.977 & 0.018 & 0.006 \\
				\bottomrule
			\end{tabular}%
		}
	\end{table}
	
	The concluding XAI table is Table ~\ref{tab:xai_evidence_summary}. Since both feature families are almost equal for the majority of models, it demonstrates that the primary distinction between explanation approaches is not whether they prioritize speed or occupancy. Temporal differences are more significant. With recent shares of 0.987 and 0.967, TreeSHAP explanations for Random Forest and XGBoost are heavily concentrated in the last 0–15 minutes before.
	
	\vspace{6pt}
	
	LIME, on the other hand, generates spread-out explanations. The recent share for Random Forest drops from 0.987 with TreeSHAP to 0.712 with LIME, although the older 30- to 60-minute timeframe shows 0.195 of the attribution. The difference is more pronounced for XGBoost: While LIME provides 0.539 and allocates 0.305 to the earlier 30- to 60-minute window, TreeSHAP provides a recent share of 0.967. This implies that TreeSHAP tends to highlight sharper recent evidence in the tree models, but LIME often catches a broader local trend across the 60-minute lookback window.
	
	\vspace{6pt}
	
	For the neural models, MLP explanations are less recent-focused than the tree-model TreeSHAP explanations, while the LSTM sequence explanation is strongly recent-focused with a recent share of 0.977. This supports the main XAI conclusion: the value of the explanations is not simply identifying speed or occupancy, but showing how recent, concentrated, and method-dependent the alarm evidence is.
	
	\vspace{6pt}
	
	\section{Comparative XAI Analysis}
	
	\subsection{Cross-Method and Cross-Model Agreement}
	
	Instead of just comparing model scores, the work compares explanations. As opposed to using raw attribution values, agreement metrics are used to compare SHAP and LIME. This is required because different mathematical ideas are used by SHAP and LIME to provide explanations. The comparison seeks to determine if the two approaches provide the same operational narrative:
	
	\begin{itemize}
		\item Do they identify the same top features?
		\item Do they focus on the same temporal region?
		\item Do they both emphasize recent traffic behavior?
	\end{itemize}
	
	Cross-model agreement measures whether different model families explain shared alarm cases using similar temporal traffic evidence. This is different from simply checking whether multiple models alarmed. Model alarm agreement asks whether models raised alarms together; explanation agreement asks whether their XAI evidence was similar.
	

	
	\begin{figure}[H]
		\centering
		\safeincludegraphics[0.65\linewidth]{cross_model_agreement_heatmap.png}
		\caption{Cross-model explanation agreement heatmap. Higher agreement means different model families explain shared alarm cases using similar temporal traffic evidence.}
		\label{fig:cross_model_heatmap}
	\end{figure}
	
	
	\begin{figure}[H]
		\centering
		\safeincludegraphics[0.9\linewidth]{cross_model_agreement_level_summary.png}
		\caption{Cross-model explanation agreement levels. The chart summarizes how often model pairs produce consistent, partially consistent, or conflicting explanations.}
		\label{fig:cross_model_levels}
	\end{figure}
	
	Figure~\ref{fig:cross_model_levels} summarizes how stable the explanation story is across model families. Agreement levels were assigned using a score based on top-5 feature overlap and contribution-direction agreement. Scores above 0.70 are treated as consistent, scores between 0.40 and 0.70 as partially consistent, and scores below 0.40 as conflicting.
	
	\section{Operator-Facing Alarm Evidence}
	
	The main value of XAI in this project is not that it says speed or occupancy mattered. Since these are the only input features, that conclusion is expected. The real value is that it turns a raw alarm into an evidence profile.
	
	An alarm evidence profile can show:
	
	\begin{itemize}
		\item whether the alarm was barely above threshold or strongly above threshold,
		\item whether one model alarmed or several models agreed,
		\item whether the explanation focused on the last 5--15 minutes or older history,
		\item whether the explanation was concentrated in a few features or spread diffusely,
		\item whether SHAP and LIME gave similar explanations,
		\item whether different models reacted to the same temporal traffic pattern.
	\end{itemize}
	
	\begin{figure}[H]
		\centering
		\safeincludegraphics[0.92\linewidth]{strong_vs_borderline_alarm_profiles.png}
		\caption{Strong versus borderline alarm profiles.}
		\label{fig:strong_borderline_profiles}
	\end{figure}
	
	Two LSTM-based operator-facing alarm evidence cards are shown in Figure ~\ref{fig:strong_borderline_profiles}. The cards are intended to convert XAI evidence and model output into a useful alarm review priority. The model, stream, anomaly score, threshold, score margin, number of alarming models, recent-evidence share, explanation concentration, and top explanatory traffic attributes are all summarized on each card.
	
	The upper example shows a strong review-worthy alarm. The LSTM score is far above the threshold, with a margin of \(0.613\), and all four models alarmed on the same case. The explanation evidence is also highly recent and concentrated, with a recent-evidence share of \(1.00\) and concentration of \(0.98\). The top evidence comes from occupancy and speed values in the last 5--10 minutes before the representative alarm timestamp. This suggests that the alarm is strongly supported by recent traffic context and should be prioritized for operator review.
	
	The lower example shows a borderline lower-priority alarm. The LSTM score is equal to the threshold, giving a margin of \(0.000\), and only one out of four models alarmed. Although the explanation is still mostly recent, the alarm is weaker because it is close to the decision threshold and has low model agreement. Therefore, the card keeps the alarm visible, but suggests lower priority unless external traffic context, such as cameras, maps, or sensor readings, confirms worsening conditions.
	
	In both examples, the plots on the right show the speed and occupancy around the alarm. The shaded region marks the alarm episode, and the dashed vertical line marks the representative row used for the XAI explanation. The figure demonstrates that the alarm support layer does not label alarms as true or false in real time. Instead, it provides structured evidence that helps the operator decide which alarms deserve immediate attention.
	
	\section{Discussion}
	
	\subsection{Operational Value}
	
	The proposed alarm support layer helps operator judgment under uncertainty. It does not replace the operator and does not claim to infer physical causality. Instead, it provides structured evidence about the model alarm.
	
	A strong, recent, concentrated, multi-model alarm can be prioritized for quick review. A borderline alarm with small score margin, weak model agreement, and diffuse explanation can remain visible but be treated as lower priority.
	
	The XAI results show that different models and methods may explain alarms differently. TreeSHAP explanations for tree models often emphasize very recent lagged features, while LIME can sometimes highlight a longer developing pattern. This is not a contradiction; it is part of the contribution. The alarm support layer can indicate when explanations are consistent, partially consistent, or conflicting, which helps avoid over-trusting a single explanation plot.
	
	\section{Limitations}
	
	The main limitations of this project are:
	
	\begin{itemize}
		\item \textbf{Narrow input feature space.} The models use only smoothed speed and
		occupancy, so XAI cannot surface evidence from variables the model never received.
		Operators should not interpret the absence of a feature explanation as evidence
		that the feature is irrelevant to the real-world incident.
		
		
		\item \textbf{LIME not available for LSTM.} The LSTM's sequential input structure
		makes standard tabular perturbation unreliable, so LIME was not applied. Cross-method
		agreement findings therefore apply only to the three flattened-input models.
		
		\item \textbf{Episode-level explanation only.} XAI was computed for one
		representative row per alarm episode. This does not capture how evidence evolves
		across the episode, so shifts in feature importance within an episode are not
		visible in the current analysis.
		
	
	\end{itemize}
	
	\section{Conclusion}
	
	This project presented an XAI-based alarm support layer for high-recall traffic incident detection. Using prepared incident-window data, four anomaly-detection model families were evaluated on smoothed speed and occupancy measurements at both timestamp and incident level.
	
	The results show that the main value of XAI in this setting is not just identifying whether speed or occupancy influenced a prediction. Since both are the only model inputs, the more relevant question is how the evidence is distributed over time. The explanation analysis showed that recent traffic conditions played a major role in many alarms, especially for TreeSHAP and the LSTM sequence attribution, while LIME sometimes reflected a wider pattern across the full lookback window.
	
	These differences are useful rather than problematic. They show that alarm interpretation can benefit from comparing explanation methods, instead of relying on one attribution view alone. Combined with score margin and model agreement, the XAI evidence provides a practical way to describe alarms as stronger, weaker, or less stable from an operator-support perspective.
	
	Overall, the project does not replace the traffic operator or prove the physical cause of an incident. Its contribution is to add interpretable evidence to binary alarms, making them easier to review under uncertainty.
	
	
	\begin{thebibliography}{9}
		
		\bibitem{abujudom2025}
		K. Abu Judom, \textit{AI-Powered Traffic Anomaly Detection with Multi-Feature Integration for Smart Mobility Solutions}, Bachelor's Thesis, Budapest University of Technology and Economics, 2025.
		
		\bibitem{lime}
		M. T. Ribeiro, S. Singh, and C. Guestrin, ``Why Should I Trust You?: Explaining the Predictions of Any Classifier,'' in \textit{Proceedings of the 22nd ACM SIGKDD International Conference on Knowledge Discovery and Data Mining}, 2016.
		
		\bibitem{shap}
		S. M. Lundberg and S.-I. Lee, ``A Unified Approach to Interpreting Model Predictions,'' in \textit{Advances in Neural Information Processing Systems}, 2017.
		

		
	\end{thebibliography}
	
\end{document}